\chapter{Introduction: the object, the objective, and the discipline}
\label{ch:intro}

\section{What problem is being solved}

A horse race is, formally, the realisation of an \emph{order statistic}. For a field of $N$
runners, each horse $i$ produces a finishing time $\Tfin_i\in\real_{>0}$, and the result of the
race --- the only thing that pays a win bet --- is the identity of the minimiser,
\begin{equation}
  \text{winner} \;=\; \argmin_{i\in\{1,\dots,N\}} \Tfin_i .
  \label{eq:argmin}
\end{equation}
Everything in this document is, ultimately, a theory of the joint law of
$(\Tfin_1,\dots,\Tfin_N)$ and of the marginal it induces on the winner through
\eqref{eq:argmin}. We do not model ``the probability that horse $i$ wins'' directly with a
classifier; we model the \emph{generative process of finishing times} and read the win
probability off it as
\begin{equation}
  p_i \;=\; \Pr\!\left[\Tfin_i < \Tfin_j \ \forall j\neq i\right].
  \label{eq:winprob}
\end{equation}
This is a deliberate and consequential choice. It means that the same machinery yields the full
finishing-order distribution (hence place and exotic prices), that distributional facts about
finishing times --- in particular \emph{variance}, not only the mean --- are first-class
citizens of the model, and that every covariate must earn its place by changing the law of a
time, not by correlating with an outcome.

\section{The objective: fit the market as goodness-of-fit}

The betting market is the most efficient publicly-available aggregator of information about
\eqref{eq:winprob}. We therefore take the market's implied (de-vigged) win probabilities as the
\emph{yardstick} and pose the modelling problem as one of \emph{goodness-of-fit to the market},
measured strictly out-of-sample and walk-forward. The headline metric is the fraction of the
market's log-likelihood advantage over a uniform baseline that the model recovers,
\begin{equation}
  \%\text{ of market} \;=\;
   \frac{\mathrm{LL}_{\text{uniform}} - \mathrm{LL}_{\text{model}}}
        {\mathrm{LL}_{\text{uniform}} - \mathrm{LL}_{\text{market}}},
  \label{eq:pctmarket}
\end{equation}
where each $\mathrm{LL}$ is the mean log-probability assigned to the actual winners over a held-out
set of races, $\mathrm{LL}_{\text{uniform}}$ uses $p_i = 1/N$, and $\mathrm{LL}_{\text{market}}$
uses the de-vigged closing line. A score of $0\%$ is no better than ignorance; $100\%$ is parity
with the closing market. At project inception the model sat near $48\%$; at the time of writing it
is at approximately $65\%$ at Sha Tin (ST) and $60\%$ at Happy Valley (HV).

Three points about this objective deserve emphasis, because they shape every later decision.

\begin{principle}[Goodness-of-fit, not edge]
$\%$ of market is a \emph{goodness-of-fit} statistic. Closing the gap to the market \emph{is}
the modelling job. The model is not a search for mispricings; whether any residual gap to the
market is tradable is a separate, downstream, user-owned decision. The modelling work is pure
distributional fit.
\end{principle}

\begin{principle}[There is always a next mechanism]
A residual between model and market is treated as a \emph{missing mechanism} until proven
otherwise. The bar for declaring a residual to be genuinely unmodellable ``private information''
is high, because that same reasoning would have terminated the programme near $10\%$ skill.
Understanding \emph{what the market knows that the model does not} is the engine of progress, not
an excuse to stop.
\end{principle}

\begin{principle}[Calibration is part of the objective]
Fit is judged jointly on discrimination (\,$\%$ of market) \emph{and} calibration (the expected
calibration error, ECE). A change that sharpens discrimination while degrading calibration is not
an improvement. The current model's ECE is $\approx 0.006$, i.e.\ predicted win frequencies match
realised ones to well under a percentage point on average.
\end{principle}

\section{The discipline: first-principle regularisation}

The defining methodological commitment of this project is that the model is regularised
\emph{by physics and game theory, not by feature count or by a penalty norm}. Every input must
couple to the generative law through a named, plausible mechanism --- a term in the equation of
motion of a horse (the ``plant''), or a parameter of the jockey/trainer control policy. An
arbitrary regressor that merely correlates with winning is forbidden, however predictive it looks
in-sample, because it both overfits and breaks the chain of physical reasoning that has carried
the model this far.

This has a sharp operational form, developed fully in Chapter~\ref{ch:assembly}: the
drift--diffusion law exposes exactly \emph{three independent knobs} a feature may load:
\begin{enumerate}[label=(\roman*)]
  \item the aerobic ceiling $\vc$ (how fast the horse can sustainably go),
  \item the anaerobic reserve $\Wp$ (the finite ``kick'' it can spend above $\vc$), and
  \item the diffusion $\Ddiff$ (the race-to-race consistency / variance of its time).
\end{enumerate}
A feature is correctly \emph{structured} only when it loads the knob its mechanism actually acts
on. Using the right \emph{signal} is necessary but not sufficient; it must enter in the right
\emph{place}. A recurring lesson --- documented repeatedly in what follows --- is that a feature
which fails to convert under one structural coupling is not a dead mechanism but a wrongly-placed
one: scalar opponent-strength (Elo) failed, but the same competition axis expressed
\emph{relationally} (Chapter~\ref{ch:competition}) converts strongly; draw failed as a mean
effect on $\vc$ but converts as a variance effect on $\Ddiff$ (Chapter~\ref{ch:geometry}).

\begin{principle}[Theory leads, statistics gate]
Theory chooses \emph{which} structure to try; statistics only gate \emph{whether that particular
structure} converts out-of-sample. A dead first attempt is not a dead axis. Conversely, a
parked verdict can go stale when the baseline or the generative law beneath it changes, and must
be re-checked against the current production path before it is trusted.
\end{principle}

\section{Leak-freeness as a non-negotiable invariant}

A backtest number is worthless --- worse, actively misleading --- if any feature has seen the
future. The entire feature layer is therefore built to a single invariant: every per-horse
quantity used to predict a race on date $d$ is a function of that horse's \emph{strictly prior}
runs only. Concretely this is enforced two ways, both of which recur throughout
Chapters~\ref{ch:pop}--\ref{ch:competition}:
\begin{itemize}
  \item \textbf{Expanding means shifted by one.} A trailing trait is the cumulative sum over a
        horse's runs \emph{minus the current run}, divided by the count \emph{minus one}:
        $\bar{x}_{<t} = (\sum_{s\le t} x_s - x_t)/(n_{\le t}-1)$. The current observation never
        enters its own feature value.
  \item \textbf{Strict as-of joins.} Where a feature lives on a separate timeline (barrier
        trials, rating updates), it is joined backward with \code{allow\_exact\_matches=False},
        so a race on date $d$ sees only updates strictly before $d$.
\end{itemize}
Contemporaneous \emph{boundary conditions} that are genuinely known before the race --- the draw,
the track geometry and rail position, the field composition, the declared carried weight --- are
public pre-race information and are allowed. The discipline is enforced mechanically by an
\code{AsOfView} wrapper and an \code{assert\_no\_lookahead} check that is run on the fit/predict
pipeline before any backtest number is trusted.

\section{How to read this document}

The chapters are ordered from the generative spine outward. Chapter~\ref{ch:gen} writes down the
finishing-time law and its surrogate. Chapter~\ref{ch:energetics} recovers the latent energetics
from data. Chapter~\ref{ch:pop} builds the partially-pooled ability population. Chapters
\ref{ch:features}--\ref{ch:competition} are the feature layer: the channels that nudge a horse's
ability away from its population prior, and the principled way to combine them.
Chapter~\ref{ch:geometry} adds the arena. Chapter~\ref{ch:assembly} assembles everything into the
production drift and the predict loop. Chapter~\ref{ch:theory} is the theory of what is not yet
modelled --- the interaction games --- and is the source of the forward research queue.
Chapters~\ref{ch:method} and~\ref{ch:state} are the method and the current state. Appendices
collect the constants table and the data surface.

Throughout, the reader should keep the three knobs and the order statistic in mind: they are the
two ideas from which everything else is derived.
